\begin{lemma}
\uses{LPCP,PCP}
\label{LPCP_to_PCP_ZMod2}
\lean{LPCP_to_PCP_ZMod2}
\leanok
\[
\LPCP[\varepsilon_c,\varepsilon_s,\Sigma=\F_{2},\ell,q,r] \subseteq
\PCP\left[
\varepsilon_c,\
\varepsilon_s'=\max\{7/8,\varepsilon_s+1/100\},
\Sigma'=\F_{2},\
\ell'=2^{\ell},\
q'=3+2\lceil\log100q\rceil q,\
r'=r+(2+\lceil\log100q\rceil)\ell
\right].
\]
\end{lemma}
\begin{proof}\uses{LINEQ_LPCP, BLR.basicVerifier, BLR_basic_completeness_ZMod2, BLR_basic_soundness_ZMod2}
TODO
\end{proof}

\begin{theorem}\label{QESAT_poly_LPCP}\lean{QESAT_poly_LPCP}\uses{QESAT, LPCP}\leanok
\[
  \mathrm{QESAT}(\F_2,n) \in
  \mathrm{LPCP}[
    \varepsilon_c = 0,
    \varepsilon_s = 3/4,
    \Sigma = \F_2,
    \ell = n + n^2,
    q = 4,
    r = |x| + 2*n
  ].
\]
\end{theorem}
\begin{proof}%\uses{} tensor stuff

\noindent Let $(p_1, \dots, p_m)$ be an instance of $\text{QESAT}(\mathbb{F})$ with $n$ variables. \\
The LPCP verifier expects a proof $\pi = (a, b) \in \mathbb{F}^{n+n^2}$ and works as follows:

\vspace{1em}
\noindent \textbf{Verifier} $V^{(a,b)}(p_1, \dots, p_m)$:
\begin{enumerate}
    \item Sample $r \gets \mathbb{F}^m$ and $s, t \gets \mathbb{F}^n$.
    \item Define 
    \[
    M := \begin{bmatrix} \text{coeff}(p_1) \\ \vdots \\ \text{coeff}(p_m) \end{bmatrix} \in \mathbb{F}^{m \times (n+n^2)} 
    \quad \text{and} \quad 
    c := \begin{bmatrix} -\text{const}(p_1) \\ \vdots \\ -\text{const}(p_m) \end{bmatrix} \in \mathbb{F}^m.
    \]
    \textit{(Note: $\text{coeff}(p_i) :=$ non-constant coeffs of $p_i$, and $\text{const}(p_i) :=$ constant coeff of $p_i$)}
    
    \item \textit{(linear equation test)} $\rightarrow$ Query $(a,b)$ at $M^T r$ and check that $\langle a \parallel b, M^T r \rangle = \langle c, r \rangle$.
    
    \item \textit{(tensor test)} $\rightarrow$ Query $b$ at $\text{flat}(s \otimes t)$, $a$ at $s$ and $t$, and check that 
    \[
    \langle b, \text{flat}(s \otimes t) \rangle = \langle a, s \rangle \cdot \langle a, t \rangle.
    \]
\end{enumerate}

\vspace{1em}
\noindent \textbf{Completeness:} 
Suppose $p_1(a) = \dots = p_m(a) = 0$ and set $b := \text{flat}(a \otimes a)$. \\
Then 
\[
\Pr_{s,t} \left[ \langle b, \text{flat}(s \otimes t) \rangle = \langle a, s \rangle \cdot \langle a, t \rangle \right] = 1
\]
and 
\[
\Pr_r \left[ \langle a \parallel b, M^T r \rangle = \langle c, r \rangle \right] = 1 \quad \left(\text{since } M \begin{bmatrix} a \\ \text{flat}(a \otimes a) \end{bmatrix} = c \right).
\]

\vspace{1em}
\noindent \textbf{Soundness:}
Suppose $\forall a \in \mathbb{F}^n \ \exists i \in [m] \text{ s.t. } p_i(a) \neq 0$. Fix any $\pi = (a, b)$. \\
Either:
\begin{enumerate}
    \item[(i)] $b \neq \text{flat}(a \otimes a) \implies$ tensor test passes w.p. $\leq \frac{2|\mathbb{F}|-1}{|\mathbb{F}|^2}$
    \item[OR (ii)] $b = \text{flat}(a \otimes a)$ and $M \begin{bmatrix} a \\ b \end{bmatrix} \neq c \implies$ linear equation test passes w.p. $\leq \frac{1}{|\mathbb{F}|}$
\end{enumerate}

\end{proof}
